\documentclass[11pt,oneside]{article}

\usepackage[a4paper,margin=27mm,headheight=15pt]{geometry}
\usepackage[T1]{fontenc}
\usepackage[utf8]{inputenc}
\IfFileExists{libertinus.sty}{\usepackage{libertinus}}{\usepackage{lmodern}}
\usepackage{microtype}
\usepackage{amsmath,amssymb,amsthm,mathtools,mathrsfs}
\usepackage{bm}
\usepackage{booktabs,longtable,array,tabularx}
\usepackage{graphicx}
\usepackage{xcolor}
\usepackage{enumitem}
\usepackage{fancyhdr}
\usepackage{titlesec}
\usepackage{xurl}
\usepackage{hyperref}
\usepackage[nameinlink,noabbrev]{cleveref}

\definecolor{linkblue}{RGB}{25,75,135}
\definecolor{shadegray}{RGB}{246,247,249}
\definecolor{rulegray}{RGB}{195,201,208}
\definecolor{provedgreen}{RGB}{27,111,72}
\definecolor{cautionorange}{RGB}{155,88,18}
\hypersetup{
  colorlinks=true,
  linkcolor=linkblue,
  citecolor=linkblue,
  urlcolor=linkblue,
  pdftitle={Gamma Duality, Total Positivity, and Lambert-Localized Moments in the Fabius--Rvachev System},
  pdfsubject={New frontier results for the Fabius function, Rvachev up-function, Thue--Morse sequence, q-binomial coefficients, sinc products, and Lambert W},
  pdfkeywords={Fabius function, Rvachev up-function, Thue-Morse, q-binomial, multiple zeta value, sinc product, generalized gamma convolution, Lambert W, mod-Gamma convergence}
}

\pagestyle{fancy}
\fancyhf{}
\fancyhead[L]{\small New Fabius--Rvachev frontiers}
\fancyhead[R]{\small\nouppercase{\leftmark}}
\fancyfoot[C]{\thepage}
\renewcommand{\headrulewidth}{0.4pt}

\titleformat{\section}{\Large\bfseries}{\thesection}{0.7em}{}
\titleformat{\subsection}{\large\bfseries}{\thesubsection}{0.7em}{}
\titleformat{\subsubsection}{\normalsize\bfseries}{\thesubsubsection}{0.7em}{}
\setcounter{secnumdepth}{3}
\setcounter{tocdepth}{2}
\setlist[itemize]{leftmargin=2em,itemsep=0.25em,topsep=0.4em}
\setlist[enumerate]{leftmargin=2.3em,itemsep=0.25em,topsep=0.4em}
\setlength{\emergencystretch}{2.5em}

\newtheorem{theorem}{Theorem}[section]
\newtheorem{proposition}[theorem]{Proposition}
\newtheorem{lemma}[theorem]{Lemma}
\newtheorem{corollary}[theorem]{Corollary}
\newtheorem{conjecture}[theorem]{Conjecture}
\theoremstyle{definition}
\newtheorem{definition}[theorem]{Definition}
\newtheorem{algorithm}[theorem]{Algorithm}
\newtheorem{example}[theorem]{Example}
\theoremstyle{remark}
\newtheorem{remark}[theorem]{Remark}
\newtheorem{warning}[theorem]{Warning}

\newcommand{\R}{\mathbb R}
\newcommand{\C}{\mathbb C}
\newcommand{\Q}{\mathbb Q}
\newcommand{\N}{\mathbb N}
\newcommand{\Z}{\mathbb Z}
\newcommand{\Up}{\operatorname{up}}
\newcommand{\sinc}{\operatorname{sinc}}
\newcommand{\e}{\mathrm e}
\newcommand{\ii}{\mathrm i}
\newcommand{\dd}{\,\mathrm d}
\newcommand{\Var}{\operatorname{Var}}
\newcommand{\Cov}{\operatorname{Cov}}
\newcommand{\E}{\mathbb E}
\newcommand{\Prob}{\mathbb P}
\newcommand{\bigO}{\mathcal O}
\newcommand{\smallo}{o}
\newcommand{\supp}{\operatorname{supp}}
\newcommand{\sgn}{\operatorname{sgn}}
\newcommand{\vtwo}{v_2}
\newcommand{\Li}{\operatorname{Li}}
\newcommand{\diag}{\operatorname{diag}}
\newcommand{\Law}{\operatorname{Law}}
\newcommand{\Cum}{\operatorname{Cum}}
\newcommand{\dist}{\mathrel{\stackrel{\mathrm d}=}}
\newcommand{\qbinom}[3]{\genfrac{[}{]}{0pt}{}{#1}{#2}_{#3}}
\newcommand{\repo}{\nolinkurl{Analysis/FabiusFunction}}
\newcommand{\statusproved}{\textcolor{provedgreen}{\textbf{proved here}}}
\newcommand{\statusfrontier}{\textcolor{cautionorange}{\textbf{frontier proposal}}}

\newenvironment{proofidea}{\par\smallskip\noindent\textit{Proof architecture.}\ }{\hfill$\square$\par\smallskip}
\newenvironment{boxedremark}{%
  \par\smallskip\noindent\begin{center}\begin{minipage}{0.94\textwidth}%
  \color{black}\hrule\vspace{0.6em}\small
}{%
  \vspace{0.6em}\hrule\end{minipage}\end{center}\smallskip
}

\title{\bfseries Gamma Duality, Total Positivity, and\\
Lambert-Localized Moments in the Fabius--Rvachev System\\[0.45em]
\large $q$-binomial multiple-zeta expansions, spectral subordinators,\\
phase-resolved endpoint tilting, and large-order transfer}
\author{Frontier research report for Vladimir Reshetnikov's \texttt{ProveIt} project}
\date{Repository snapshot: 27 August 2026\\
\small commit \nolinkurl{158dd2e12d7412142d79a53d91d8fcbdfdfb5d97}}

\begin{document}
\maketitle

\begin{abstract}
The Fabius--Rvachev corpus already contains a remarkably complete triangle of
representations: a dyadic random series, exact $q$-binomial arithmetic, and an
infinite product of sinc functions whose negative-Laplace endpoint asymptotic
contains a one-periodic fluctuation and the branch $W_{-1}$ of Lambert's
function.  This report develops a new layer connecting those three sides.
All claims marked ``proved here'' are supplied with proofs in the present
manuscript; they are not claimed to have been formalized in Lean yet.

The first result expands every even moment simultaneously in Gaussian
$q$-binomial coefficients and symmetrized multiple zeta values.  The Wick-rotated
sinc determinant is then shown to lie in the type-I Laguerre--P\'olya class.
Consequently its coefficient sequence is P\'olya-frequency of infinite order,
and every shifted Jensen polynomial is real-rooted; this gives an infinite
hierarchy of sharp moment inequalities.  The same determinant proves that the
Rvachev law is strictly sub-Gaussian with the optimal variance proxy, although
its real-axis zeros rule out infinite divisibility.

The principal structural theorem identifies the reciprocal Wick-rotated
product as a Laplace transform.  Every positive power is the transform of an
explicit generalized gamma convolution with discrete Thorin measure
$\tau\sum_{m\ge1}(1+v_2(m))\delta_{m^2}$.  Its cumulants are
$\tau(r-1)!\zeta(2r)/(1-4^{-r})$, it satisfies an exact $q=1/4$ perpetuity, and
the parity of its cumulative Thorin mass is precisely the Thue--Morse sign of
the Fourier lobes.  The associated L\'evy heat kernel admits not only the usual
Mellin expansion with log-periodic oscillation, but an exact Jacobi-theta
completion whose error begins at
$-2\sqrt{\pi/x}\exp(-4\pi^2/x)$.

Finally, exponential tilting of the endpoint law leaves the binary coordinates
independent and turns them into truncated exponentials.  This yields a
phase-resolved mod-Gamma limit, a central limit theorem with Berry--Esseen error
$\bigO((\log s)^{-1/2})$, and a rigorous Lambert-$W_{-1}$ localization law.  A
universal endpoint-transfer theorem then converts the exact negative-Laplace
expansion into all-orders asymptotics for ordinary moments and centered even
moments.  In particular it supplies the previously missing bridge from the
exact $q$-binomial/multiple-zeta coefficients of the sinc determinant to the
Lambert-$W$ endpoint saddle.
\end{abstract}

\tableofcontents
\newpage

\section{Corpus audit, baseline, and novelty discipline}

\subsection{Scope of the audit}

The audited source tree is \repo/\texttt{docs} at the commit printed on the
title page.  It contains fifty-seven standalone \LaTeX{} manuscripts when
historical revisions, source transcriptions of reference papers, and the live
frontier drafts are counted separately.  They fall into five overlapping
families:
\begin{enumerate}
\item the current primary exposition, Lean walkthrough, glossary, and
      non-elementarity study;
\item archived primary expositions, missing-parts supplements, and
      $q$-calculus introductions;
\item exact-dyadic, Thue--Morse, repeated-integration, inverse-function, and
      small-argument studies;
\item the consolidated non-formalized frontier and its sequence of frontier
      reports;
\item the multi-wave Fourier-decay campaign and the transcribed reference
      literature.
\end{enumerate}
The revisions were not treated as independent evidence: later corrections
supersede earlier drafts.  In particular, the current Fourier audit already
contains the Thue--Morse lobe-sign theorem, exact zero orders, the
renormalization identity, and the global sharp decay envelope.  Those are
baseline results here, not rediscoveries.

\subsection{What was already present}

The corpus supplies, among much else, the following ingredients used below:
\begin{itemize}
\item the random-series realization of the bounded Fabius distribution and of
      Rvachev's even density;
\item the sinc product, its integer zero multiplicities
      $a_m=1+v_2(m)$, and its dyadic renormalization;
\item exact moment recurrences, rationality, finite convolution models, and
      dyadic values expressed through Gaussian binomial coefficients;
\item spectral-zeta and determinant viewpoints for the Wick-rotated product;
\item the exact negative-Laplace decomposition
      \[
      \log L(s)=-\frac{(\log s)^2}{2\log2}+\frac{\log s}{2}
        +R(\log_2s)+E(s),
      \]
      where $R$ is one-periodic and $E(s)$ is positive and exponentially
      small;
\item Lambert-$W_{-1}$ descriptions of the small-argument endpoint saddle;
\item exact Thue--Morse signs of the Fourier lobes and sharp real-axis decay.
\end{itemize}
The goal was therefore not to enlarge the collection by cosmetic reformulation,
but to find implications which genuinely cross the established boundaries
between $q$-arithmetic, entire-function theory, probability, and endpoint
asymptotics.

\subsection{Status ledger}

\begin{longtable}{@{}p{0.07\textwidth}p{0.66\textwidth}p{0.19\textwidth}@{}}
\toprule
ID & Result & Status\\
\midrule
\endhead
A & Gaussian $q$-binomial/occupancy expansion of every even moment in
    symmetrized multiple zeta values, with a partition-lattice reduction to
    ordinary even zeta values. & \statusproved\\
B & Type-I Laguerre--P\'olya classification, PF$_\infty$ coefficients,
    hyperbolic shifted Jensen polynomials, and the resulting ultra-log-concave
    moment hierarchy. & \statusproved\\
C & Strict sub-Gaussianity with the exact variance proxy, together with
    non-infinite-divisibility. & \statusproved\\
D & Explicit generalized gamma-convolution dual, discrete Thorin measure,
    cumulant duality, and an exact $q=1/4$ perpetuity. & \statusproved\\
E & Thue--Morse lobe sign as parity of cumulative Thorin mass. & \statusproved\\
F & Exact Jacobi-theta completion of the L\'evy heat kernel, including its
    first exponentially flat correction. & \statusproved\\
G & Phase-resolved mod-Gamma convergence of the endpoint tilt and a
    Berry--Esseen central limit theorem. & \statusproved\\
H & Universal endpoint-transfer expansion, applied to ordinary moments,
    centered moments, and the large-order $q$/MZV coefficients. & \statusproved\\
I & Phase-resolved local limit and full Edgeworth expansion under the endpoint
    tilt. & \statusfrontier\\
\bottomrule
\end{longtable}

\begin{boxedremark}
The word ``new'' in this report means new to the audited corpus after targeted
searches for the exact structures used below.  It is not a claim of priority
over every article in the mathematical literature.  The general theories of
multiple zeta values, P\'olya-frequency sequences, generalized gamma
convolutions, Mellin harmonic sums, and Lambert $W$ are classical; what is new
here is their exact synthesis for the Fabius--Rvachev product and the resulting
cross-representation theorems.
\end{boxedremark}

\section{Normalization and the master determinant}

\subsection{The random series}

Let $(V_k)_{k\ge1}$ be independent uniform random variables on $[0,1]$, and let
$U_k=2V_k-1$, uniform on $[-1,1]$.  Define
\begin{equation}\label{eq:X-Y-def}
 X=\sum_{k\ge1}2^{-k}V_k,\qquad
 Y=2X-1=\sum_{k\ge1}2^{-k}U_k.
\end{equation}
Then $X\in[0,1]$, $Y\in[-1,1]$, and
\begin{equation}\label{eq:symmetry}
 X\dist1-X,\qquad Y\dist-Y.
\end{equation}
The distribution function of $X$ is the bounded Fabius function, and the density
of $Y$ is Rvachev's $\Up$-function.  We write
\begin{equation}\label{eq:moments-def}
 m_n=\E X^n,\qquad \mu_{2n}=\E Y^{2n}.
\end{equation}

\subsection{Fourier and Wick-rotated products}

With $\sinc z=\sin z/z$ and the Fourier convention
$\Phi(z)=\E\e^{-2\pi\ii zY}$,
\begin{equation}\label{eq:Phi-products}
 \Phi(z)=\prod_{j=0}^{\infty}\sinc\!\left(\frac{\pi z}{2^j}\right)
 =\prod_{m=1}^{\infty}\left(1-\frac{z^2}{m^2}\right)^{a_m},
 \qquad a_m=1+v_2(m).
\end{equation}
The first product gives the exact renormalization
\begin{equation}\label{eq:Phi-renorm}
 \Phi(z)=\sinc(\pi z)\Phi(z/2),
\end{equation}
whereas the second shows that the zero at a nonzero integer $m$ has order
$a_m$.

The central entire function of this report is the Wick rotation
\begin{equation}\label{eq:A-def}
 \mathcal A(w)=\Phi(\ii\sqrt w)
 =\prod_{m\ge1}\left(1+\frac{w}{m^2}\right)^{a_m}
 =\prod_{j\ge0}\prod_{n\ge1}\left(1+\frac{4^{-j}w}{n^2}\right).
\end{equation}
It is independent of the branch of $\sqrt w$ because it is entire in $w$.
Set
\begin{equation}\label{eq:C-def}
 \mathcal A(w)=\sum_{r\ge0}C_rw^r.
\end{equation}
Since
\begin{equation}\label{eq:mgf-A}
 \E\e^{tY}=\mathcal A\!\left(\frac{t^2}{4\pi^2}\right),
\end{equation}
we have the exact coefficient--moment correspondence
\begin{equation}\label{eq:C-moment}
 C_r=\frac{(2\pi)^{2r}}{(2r)!}\,\mu_{2r}.
\end{equation}

\subsection{The spectral power sums}

The logarithm of \cref{eq:A-def} is
\begin{equation}\label{eq:logA-power}
 \log\mathcal A(w)
 =\sum_{r\ge1}\frac{(-1)^{r+1}}r p_rw^r,
 \qquad
 p_r=\sum_{m\ge1}\frac{a_m}{m^{2r}}
 =\frac{\zeta(2r)}{1-4^{-r}}.
\end{equation}
Consequently
\begin{equation}\label{eq:C-recurrence}
 rC_r=\sum_{j=1}^{r}(-1)^{j+1}p_jC_{r-j},\qquad C_0=1.
\end{equation}
This Newton recurrence is already enough to prove rationality of all
$\mu_{2r}$ because every $\zeta(2j)$ is a rational multiple of $\pi^{2j}$.
The next section gives a different, genuinely $q$-binomial route which keeps
the dyadic layers visible.

\section{A $q$-binomial and multiple-zeta formula for every even moment}

Throughout this section
\begin{equation}\label{eq:q-def}
 q=\frac14.
\end{equation}

\subsection{Finite depth and one-site occupancy}

For $J\ge1$ truncate the dyadic layers:
\begin{equation}\label{eq:AJ-def}
 \mathcal A_J(w)=\prod_{j=0}^{J-1}\prod_{n\ge1}
 \left(1+\frac{q^jw}{n^2}\right)
 =\sum_{r\ge0}C_r^{(J)}w^r.
\end{equation}
At a fixed spectral site $n$, the finite $q$-binomial theorem gives
\begin{equation}\label{eq:finite-qbinom}
 \prod_{j=0}^{J-1}(1+q^jx)
 =\sum_{r=0}^{J}b_{J,r}x^r,
 \qquad
 b_{J,r}=q^{r(r-1)/2}\qbinom{J}{r}{q}.
\end{equation}
Thus $r$ records the number of dyadic layers occupying the same integer
frequency $n$.  In the infinite-depth limit,
\begin{equation}\label{eq:b-infty}
 b_r=\lim_{J\to\infty}b_{J,r}
 =\frac{q^{r(r-1)/2}}{(q;q)_r},
 \qquad (q;q)_r=\prod_{j=1}^{r}(1-q^j).
\end{equation}
For $q=1/4$,
\begin{equation}\label{eq:b-first}
 b_1=\frac43,\qquad b_2=\frac{16}{45},\qquad
 b_3=\frac{64}{2835}.
\end{equation}

\subsection{Distinct-index zeta sums}

Let $\lambda=(\lambda_1,\dots,\lambda_\ell)$ be a partition of $r$, and let
$m_j(\lambda)$ denote the multiplicity of the part $j$.  Define the
symmetrized distinct-index zeta sum
\begin{equation}\label{eq:Zlambda-def}
 Z_\lambda=
 \sum_{\substack{n_1,\dots,n_\ell\ge1\\ n_i\ne n_j\ (i\ne j)}}
 \frac1{n_1^{2\lambda_1}\cdots n_\ell^{2\lambda_\ell}}.
\end{equation}
This notation keeps the assignment of distinct spectral sites to the parts of
$\lambda$ explicit.

\begin{theorem}[$q$-binomial/MZV moment formula]\label{thm:q-mzv}
For every $r\ge0$ and $J\ge1$,
\begin{equation}\label{eq:CJ-partition}
 \boxed{
 C_r^{(J)}=
 \sum_{\lambda\vdash r}
 \frac{Z_\lambda}{\prod_{j\ge1}m_j(\lambda)!}
 \prod_{i=1}^{\ell(\lambda)}b_{J,\lambda_i}.}
\end{equation}
Passing to infinite depth gives
\begin{equation}\label{eq:C-partition}
 \boxed{
 C_r=
 \sum_{\lambda\vdash r}
 \frac{Z_\lambda}{\prod_{j\ge1}m_j(\lambda)!}
 \prod_{i=1}^{\ell(\lambda)}
 \frac{q^{\lambda_i(\lambda_i-1)/2}}{(q;q)_{\lambda_i}}.}
\end{equation}
Moreover
\begin{equation}\label{eq:Z-MZV}
 Z_\lambda=\sum_{\sigma\in S_\ell}
 \zeta(2\lambda_{\sigma(1)},\dots,2\lambda_{\sigma(\ell)}),
\end{equation}
where
$\zeta(s_1,\dots,s_\ell)=\sum_{n_1>\cdots>n_\ell\ge1}
 n_1^{-s_1}\cdots n_\ell^{-s_\ell}$.
\end{theorem}

\begin{proof}
Apply \cref{eq:finite-qbinom} at each spectral site $n$:
\[
 \mathcal A_J(w)=\prod_{n\ge1}
 \left(\sum_{r=0}^{J}b_{J,r}\frac{w^r}{n^{2r}}\right).
\]
A monomial of total degree $r$ is specified by a finite set of distinct sites
$n_1,\dots,n_\ell$ and positive occupancies
$\lambda_1+\cdots+\lambda_\ell=r$.  Sorting the occupancies gives a partition
$\lambda$.  The sum over ordered, pairwise distinct sites is $Z_\lambda$;
permuting sites carrying equal occupancies does not change the monomial, hence
the divisor $\prod_jm_j(\lambda)!$.  This proves
\cref{eq:CJ-partition}.  Since the products converge normally on compact
subsets and $b_{J,r}\to b_r$, coefficientwise convergence proves
\cref{eq:C-partition}.

Finally, every pairwise-distinct $\ell$-tuple has a unique strict ordering.
Summing over the $\ell!$ possible orderings of the labelled coordinates gives
\cref{eq:Z-MZV}; repeated parts correctly produce repeated MZVs, matching the
labelled sum in \cref{eq:Zlambda-def}.
\end{proof}

\subsection{Reduction to ordinary zeta values}

The MZV form is conceptually natural, but in the present fully symmetrized
combination it collapses to ordinary zeta values by M\"obius inversion on the
partition lattice.

\begin{theorem}[Partition-lattice reduction]\label{thm:partition-lattice}
Let $\Pi_\ell$ be the lattice of set partitions of $\{1,\dots,\ell\}$.  Then
\begin{equation}\label{eq:Z-partition-lattice}
 \boxed{
 Z_\lambda=
 \sum_{\pi\in\Pi_\ell}
 \left(\prod_{B\in\pi}(-1)^{|B|-1}(|B|-1)!\right)
 \prod_{B\in\pi}
 \zeta\!\left(2\sum_{i\in B}\lambda_i\right).}
\end{equation}
Consequently \cref{eq:C-partition} is a finite expression involving only
$q$-Pochhammer factors and ordinary even zeta values.
\end{theorem}

\begin{proof}
Start with the unrestricted sum over $(n_1,\dots,n_\ell)$.  Coincidences among
indices are encoded by set partitions: a block $B$ means that all $n_i$ with
$i\in B$ coincide.  M\"obius inversion from the full coincidence lattice to
the discrete partition has M\"obius factor
$(-1)^{|B|-1}(|B|-1)!$ on each block.  The sum attached to a block is
$\zeta(2\sum_{i\in B}\lambda_i)$.  Multiplication over blocks yields
\cref{eq:Z-partition-lattice}.
\end{proof}

\begin{corollary}[A second rationality proof]\label{cor:q-mzv-rational}
Every centered even moment $\mu_{2r}$ is rational.  More precisely,
$C_r/\pi^{2r}\in\Q$, and \cref{eq:C-moment} gives $\mu_{2r}\in\Q$.
\end{corollary}

\begin{remark}
The usual Newton recurrence \cref{eq:C-recurrence} forgets where repeated
spectral integers came from.  Formula \cref{eq:C-partition} retains that
information: the Gaussian $q$-binomial factor is an occupancy weight across
dyadic layers, while the MZV records the strict ordering of distinct integer
frequencies.  This is the precise meeting point of the $q$-binomial and sinc
representations.
\end{remark}

\subsection{First examples}

Using either \cref{eq:C-partition} or \cref{eq:C-recurrence},
\begin{align}
 C_1&=\frac{2\pi^2}{9},
 &\mu_2&=\frac19,\label{eq:first-moments-a}\\
 C_2&=\frac{38\pi^4}{2025},
 &\mu_4&=\frac{19}{675},\\
 C_3&=\frac{2332\pi^6}{2679075},
 &\mu_6&=\frac{583}{59535},\\
 C_4&=\frac{265618\pi^8}{10247461875},
 &\mu_8&=\frac{132809}{32531625}.\label{eq:first-moments-b}
\end{align}
For example, at degree two the partitions $(2)$ and $(1,1)$ give
\begin{equation}\label{eq:C2-MZV}
 C_2=b_2\zeta(4)+b_1^2\zeta(2,2).
\end{equation}
The identity $2\zeta(2,2)=\zeta(2)^2-\zeta(4)$ immediately recovers the
Newton form $(p_1^2-p_2)/2$.

\section{Laguerre--P\'olya structure and an infinite hierarchy of moment inequalities}

\subsection{Classification of the Wick-rotated determinant}

\begin{theorem}[Type-I Laguerre--P\'olya classification]\label{thm:LP}
The entire function $\mathcal A$ belongs to the type-I Laguerre--P\'olya class.
It has genus zero, order $1/2$, and only real nonpositive zeros, namely
$-m^2$ with multiplicity $a_m=1+v_2(m)$.
\end{theorem}

\begin{proof}
The spectral sum
\[
 \sum_{m\ge1}\frac{a_m}{m^2}=\frac{\zeta(2)}{1-4^{-1}}
 =\frac{2\pi^2}{9}
\]
converges.  Therefore the genus-zero product in \cref{eq:A-def} converges
normally on compact sets.  Every finite truncation is a polynomial with all
zeros real and nonpositive; their locally uniform limit is $\mathcal A$, so
$\mathcal A$ is type-I Laguerre--P\'olya.  The exponent of convergence of the
zeros is $1/2$, because
\[
 \sum_{m\ge1}\frac{a_m}{(m^2)^\rho}
 =\frac{\zeta(2\rho)}{1-4^{-\rho}}
\]
converges exactly for $\rho>1/2$.  The support bound $|Y|\le1$ and
\cref{eq:mgf-A} show that $\mathcal A(w)$ has growth
$\exp(\bigO(\sqrt{|w|}))$, completing the order statement.
\end{proof}

\subsection{P\'olya frequency and Jensen hyperbolicity}

Write
\begin{equation}\label{eq:gamma-def}
 \mathcal A(w)=\sum_{r\ge0}\gamma_r\frac{w^r}{r!},
 \qquad
 \gamma_r=r!C_r
 =\frac{r!(2\pi)^{2r}}{(2r)!}\mu_{2r}.
\end{equation}

\begin{theorem}[PF$_\infty$ and Jensen polynomials]\label{thm:PF-Jensen}
The ordinary coefficient sequence $(C_r)_{r\ge0}$ is a P\'olya-frequency
sequence of infinite order.  Equivalently, every finite minor of the Toeplitz
matrix
\begin{equation}\label{eq:Toeplitz}
 T=(C_{j-i})_{i,j\ge0},\qquad C_k=0\quad(k<0),
\end{equation}
is nonnegative.

Furthermore, for every $n,d\ge0$, the shifted Jensen polynomial
\begin{equation}\label{eq:Jensen}
 J_n^{(d)}(x)=\sum_{j=0}^{d}\binom dj\gamma_{n+j}x^j
\end{equation}
has only real nonpositive zeros.
\end{theorem}

\begin{proof}
The Aissen--Schoenberg--Whitney characterization says that the generating
function of a PF$_\infty$ sequence has the form
\[
 Cz^me^{\sigma z}\frac{\prod_i(1+\alpha_i z)}
 {\prod_j(1-\beta_jz)},
 \qquad \alpha_i,\beta_j,\sigma\ge0,
 \quad \sum_i\alpha_i+\sum_j\beta_j<\infty.
\]
Equation \cref{eq:A-def} is exactly this form with no denominator, no
exponential factor, and $\alpha_i=1/m^2$ repeated $a_m$ times.  This proves the
PF$_\infty$ assertion.  Jensen's characterization of the Laguerre--P\'olya
class, applied to \cref{eq:gamma-def}, gives hyperbolicity of all shifted
Jensen polynomials.  Their coefficients are positive, so their real zeros
cannot be positive.
\end{proof}

\begin{corollary}[Ultra-log-concavity of even moments]\label{cor:moment-turan}
For every $r\ge1$,
\begin{equation}\label{eq:gamma-turan}
 \gamma_r^2\ge\gamma_{r-1}\gamma_{r+1}.
\end{equation}
Equivalently,
\begin{equation}\label{eq:mu-turan}
 \boxed{
 \mu_{2r}^2\ge
 \frac{2r-1}{2r+1}\,\mu_{2r-2}\mu_{2r+2}.}
\end{equation}
Hence the sequence
\begin{equation}\label{eq:ratio-monotone}
 \frac{2\pi^2}{2r+1}\frac{\mu_{2r+2}}{\mu_{2r}}
 =\frac{\gamma_{r+1}}{\gamma_r}
\end{equation}
is nonincreasing.
\end{corollary}

\begin{proof}
The discriminant of
$J_{r-1}^{(2)}(x)=\gamma_{r-1}+2\gamma_rx+\gamma_{r+1}x^2$
is nonnegative.  Substituting \cref{eq:gamma-def} and cancelling the common
powers of $2\pi$ gives \cref{eq:mu-turan}.  The ratio statement is the usual
consequence of log-concavity.
\end{proof}

\begin{remark}[Higher inequalities]
The quadratic inequality is only the first shadow of
\cref{thm:PF-Jensen}.  Every Toeplitz minor and every discriminant or subdiscriminant
of every $J_n^{(d)}$ gives a polynomial inequality among exact Fabius moments.
This hierarchy is algorithmic and strictly stronger than checking Hankel
positivity, which follows for unrelated reasons from the fact that
$(\mu_{2r})$ is the moment sequence of $Y^2$.
\end{remark}

\section{Optimal sub-Gaussian concentration without infinite divisibility}

\begin{theorem}[Strict optimal sub-Gaussianity]\label{thm:subgaussian}
For every real $t$,
\begin{equation}\label{eq:subgaussian}
 \E\e^{tY}\le\exp\!\left(\frac{t^2}{18}\right),
\end{equation}
with strict inequality for $t\ne0$.  The variance proxy $1/9$ is optimal.
Consequently, for $u\ge0$,
\begin{equation}\label{eq:tails-YX}
 \Prob(|Y|\ge u)\le2\e^{-9u^2/2},\qquad
 \Prob\!\left(\left|X-\frac12\right|\ge u\right)
 \le2\e^{-18u^2}.
\end{equation}
\end{theorem}

\begin{proof}
From \cref{eq:mgf-A,eq:A-def},
\[
 \log\E\e^{tY}
 =\sum_{m\ge1}a_m
 \log\!\left(1+\frac{t^2}{4\pi^2m^2}\right).
\]
The strict inequality $\log(1+x)<x$ for $x>0$ gives, when $t\ne0$,
\[
 \log\E\e^{tY}<\frac{t^2}{4\pi^2}
 \sum_{m\ge1}\frac{a_m}{m^2}
 =\frac{t^2}{18}.
\]
At the origin equality holds.  Since $\Var(Y)=\mu_2=1/9$, no smaller
quadratic proxy can dominate the moment generating function near zero.
Chernoff optimization gives the two tail bounds.
\end{proof}

\begin{theorem}[The Fabius--Rvachev law is not infinitely divisible]
\label{thm:not-ID}
Neither $X$ nor $Y$ is infinitely divisible; in particular, neither is
self-decomposable.
\end{theorem}

\begin{proof}
The characteristic function of $Y$ is
$t\mapsto\Phi(t/(2\pi))$, and \cref{eq:Phi-products} shows that it vanishes at
$t=2\pi m$ for every nonzero integer $m$.  An infinitely divisible
characteristic function is the exponential of a L\'evy--Khintchine exponent
and hence never vanishes.  Translation and nonzero scaling preserve infinite
divisibility, so the assertion transfers between $X$ and $Y$.
\end{proof}

\begin{boxedremark}
There is no contradiction between \cref{thm:subgaussian,thm:not-ID}.  The law
is as concentrated as its variance permits in the sub-Gaussian sense, but it
does not sit inside a convolution semigroup.  The next section identifies a
canonical object on the other side of the sinc determinant which \emph{is}
infinitely divisible.
\end{boxedremark}

\section{The generalized gamma-convolution dual}

\subsection{Construction and Thorin measure}

For $\tau>0$, let $G_{j,n}^{(\tau)}$ be independent gamma variables with shape
$\tau$ and rate $n^2$:
\begin{equation}\label{eq:gamma-LT}
 \E\e^{-wG_{j,n}^{(\tau)}}=\left(1+\frac{w}{n^2}\right)^{-\tau}.
\end{equation}

\begin{theorem}[Spectral gamma dual]\label{thm:gamma-dual}
The nonnegative random series
\begin{equation}\label{eq:Vtau-layered}
 V_\tau=\sum_{j\ge0}\sum_{n\ge1}q^jG_{j,n}^{(\tau)},
 \qquad q=\frac14,
\end{equation}
converges almost surely and in $L^1$, with
\begin{equation}\label{eq:V-mean}
 \E V_\tau=\frac{2\tau\pi^2}{9}.
\end{equation}
Its Laplace transform is exactly
\begin{equation}\label{eq:V-LT}
 \boxed{\E\e^{-wV_\tau}=\mathcal A(w)^{-\tau}},\qquad w\ge0.
\end{equation}
Equivalently,
\begin{equation}\label{eq:V-grouped}
 V_\tau\dist\sum_{m\ge1}H_m,
 \qquad
 H_m\sim\operatorname{Gamma}(\tau a_m,\text{rate }m^2)
\end{equation}
with independent summands.  Thus $V_\tau$ is a generalized gamma convolution
with discrete Thorin measure
\begin{equation}\label{eq:Thorin}
 \boxed{U_\tau=\tau\sum_{m\ge1}a_m\,\delta_{m^2}.}
\end{equation}
\end{theorem}

\begin{proof}
The expectation of the nonnegative series is
\[
 \tau\sum_{j,n}\frac{q^j}{n^2}
 =\frac{\tau\zeta(2)}{1-q}=\frac{2\tau\pi^2}{9}<\infty,
\]
so the series is finite almost surely and converges in $L^1$.  Independence and
\cref{eq:gamma-LT} give
\[
 \E\e^{-wV_\tau}
 =\prod_{j,n}\left(1+\frac{q^jw}{n^2}\right)^{-\tau}
 =\mathcal A(w)^{-\tau}.
\]
Regrouping equal rates uses the identity that the sum of independent gamma
variables with a common rate is gamma with the sum of their shapes.  The
multiplicity of rate $m^2$ is $a_m$.  Finally, the defining Laplace exponent of
a generalized gamma convolution is
$\int\log(1+w/u)\,U(\dd u)$; equation \cref{eq:Thorin} reproduces
$\tau\log\mathcal A(w)$.
\end{proof}

\begin{corollary}[Exact cumulants and cumulant duality]\label{cor:V-cumulants}
For every $r\ge1$,
\begin{equation}\label{eq:V-cumulants}
 \boxed{
 \kappa_r(V_\tau)=\tau(r-1)!\frac{\zeta(2r)}{1-4^{-r}}.}
\end{equation}
The even cumulants of $Y$ are obtained from those of $V_1$ by
\begin{equation}\label{eq:cumulant-duality}
 \boxed{
 \kappa_{2r}(Y)=(-1)^{r+1}
 \frac{(2r)!}{r!(2\pi)^{2r}}\,\kappa_r(V_1).}
\end{equation}
\end{corollary}

\begin{proof}
A gamma variable with shape $a$ and rate $\lambda$ has $r$th cumulant
$a(r-1)!/\lambda^r$.  Summing \cref{eq:V-grouped} gives
\cref{eq:V-cumulants}.  Comparing the expansion
\[
 \log\E\e^{tY}=\log\mathcal A(t^2/(4\pi^2))
\]
with \cref{eq:logA-power} proves \cref{eq:cumulant-duality}.
\end{proof}

\subsection{An exact $q=1/4$ perpetuity}

Define the one-layer variable
\begin{equation}\label{eq:Btau}
 B_\tau=\sum_{n\ge1}G_{0,n}^{(\tau)}.
\end{equation}
Euler's product for $\sinh$ gives
\begin{equation}\label{eq:B-LT}
 \E\e^{-wB_\tau}
 =\left(\frac{\pi\sqrt w}{\sinh(\pi\sqrt w)}\right)^\tau.
\end{equation}

\begin{corollary}[Gamma perpetuity]\label{cor:perpetuity}
Let $V_\tau'$ be an independent copy of $V_\tau$.  Then
\begin{equation}\label{eq:perpetuity}
 \boxed{V_\tau\dist B_\tau+\frac14V_\tau'.}
\end{equation}
This equation has a unique law on $[0,\infty)$ with finite first moment.
\end{corollary}

\begin{proof}
Split the layer $j=0$ from \cref{eq:Vtau-layered}; after shifting
$j\mapsto j+1$, the remaining sum is $q$ times an independent copy.  Uniqueness
follows by iteration in the Wasserstein-$1$ metric, where the map
$\nu\mapsto\Law(B_\tau+qZ)$ is a contraction with factor $q$.
\end{proof}

\subsection{Thue--Morse parity as Thorin-mass parity}

Let $s_2(N)$ be the number of $1$'s in the binary expansion of $N$.
Legendre's identity $v_2(N!)=N-s_2(N)$ gives
\begin{equation}\label{eq:cumulative-a}
 \sum_{m\le N}a_m=N+v_2(N!)=2N-s_2(N).
\end{equation}

\begin{theorem}[Thue--Morse/Thorin parity law]\label{thm:Thorin-parity}
For every integer $N\ge1$,
\begin{equation}\label{eq:Thorin-parity}
 \boxed{
 (-1)^{U_1([1,N^2])}=(-1)^{s_2(N)}.}
\end{equation}
Hence, for every positive noninteger $x$,
\begin{equation}\label{eq:lobe-Thorin}
 \boxed{
 \sgn\Phi(x)=(-1)^{U_1([1,\lfloor x\rfloor^2])}.}
\end{equation}
\end{theorem}

\begin{proof}
By \cref{eq:Thorin}, $U_1([1,N^2])=\sum_{m\le N}a_m$.
Equation \cref{eq:cumulative-a} has the same parity as $s_2(N)$.  The current
Fourier sign theorem gives
$\sgn\Phi(x)=(-1)^{s_2(\lfloor x\rfloor)}$ on each open lobe; substitution
proves \cref{eq:lobe-Thorin}.
\end{proof}

\begin{remark}
The Thue--Morse sign is therefore not merely an external combinatorial label.
It is the parity of the spectral mass accumulated by the infinitely divisible
dual before the corresponding squared zero.  This interpretation links the
real-axis sign changes of a non-infinitely-divisible compactly supported law to
the Thorin measure of its reciprocal Wick rotation.
\end{remark}

\subsection{Complete Bernstein and Stieltjes consequences}

Differentiating the logarithm gives
\begin{equation}\label{eq:Stieltjes-logderivative}
 \frac{\mathcal A'(w)}{\mathcal A(w)}
 =\sum_{m\ge1}\frac{a_m}{m^2+w}
 =\int_{[1,\infty)}\frac{U_1(\dd u)}{u+w}.
\end{equation}
Thus $\mathcal A'/\mathcal A$ is a Stieltjes function,
$\log\mathcal A$ is a complete Bernstein function, and
$\mathcal A^{-\tau}$ is completely monotone for every $\tau>0$.  Standard
Stieltjes theory then supplies convergent continued fractions and alternating
Pad\'e bounds on the positive real axis.  These are certified rational bounds
for the reciprocal Wick-rotated sinc product, complementary to the product
truncations used elsewhere in the corpus.

\section{The L\'evy heat kernel and its exact theta completion}

\subsection{L\'evy density and Mellin transform}

The gamma dual has L\'evy--Khintchine representation
\begin{equation}\label{eq:Levy-Khintchine}
 -\log\E\e^{-wV_\tau}
 =\tau\int_0^\infty(1-\e^{-wx})K(x)\frac{\dd x}{x},
\end{equation}
where
\begin{equation}\label{eq:K-def}
 K(x)=\sum_{m\ge1}a_m\e^{-m^2x}
 =\sum_{j\ge0}\sum_{n\ge1}\e^{-4^jn^2x}.
\end{equation}
Its Mellin transform is
\begin{equation}\label{eq:K-Mellin}
 \int_0^\infty K(x)x^{s-1}\dd x
 =\Gamma(s)\frac{\zeta(2s)}{1-4^{-s}},
 \qquad \Re s>\frac12.
\end{equation}
The poles at
\begin{equation}\label{eq:omega-k}
 \omega_k=\frac{\pi\ii k}{\log2},\qquad k\in\Z,
\end{equation}
produce a periodic fluctuation in $\log_4(1/x)$.

Put $L_2=\log2$ and
\begin{equation}\label{eq:CK}
 C_K=\frac{\gamma}{4L_2}-\frac14-
 \frac{\log(2\pi)}{2L_2}
 =-1.36756152041694261136\ldots,
\end{equation}
\begin{equation}\label{eq:PK}
 P_K(u)=\frac1{2L_2}\sum_{k\ne0}
 \Gamma(\omega_k)\zeta(2\omega_k)\e^{2\pi\ii ku}.
\end{equation}
The series is real-valued and absolutely, exponentially convergent.
Define
\begin{equation}\label{eq:MK}
 M_K(x)=\sqrt{\frac\pi x}
 -\frac{\log(1/x)}{4L_2}+C_K
 +P_K\!\left(\log_4\frac1x\right).
\end{equation}
Mellin residue calculus gives $K(x)=M_K(x)+\bigO_N(x^N)$ for every $N>0$.
The next theorem identifies the flat remainder exactly.

\subsection{Exact completion}

Let
\begin{equation}\label{eq:theta-plus}
 \vartheta_+(t)=\sum_{n\ge1}\e^{-n^2t}.
\end{equation}

\begin{theorem}[Exact Jacobi-theta completion]\label{thm:theta-completion}
For every $x>0$,
\begin{equation}\label{eq:theta-completion}
 \boxed{
 K(x)=M_K(x)-2\sqrt{\frac\pi x}
 \sum_{r=0}^{\infty}2^r
 \vartheta_+\!\left(\frac{4^{r+1}\pi^2}{x}\right).}
\end{equation}
In particular, as $x\downarrow0$,
\begin{equation}\label{eq:K-flat-leading}
 \boxed{
 K(x)-M_K(x)=
 -2\sqrt{\frac\pi x}\e^{-4\pi^2/x}
 \left(1+\bigO(\e^{-12\pi^2/x})\right).}
\end{equation}
\end{theorem}

\begin{proof}
From the layered definition,
\begin{equation}\label{eq:K-difference}
 K(x)-K(4x)=\vartheta_+(x).
\end{equation}
The Jacobi transformation is
\begin{equation}\label{eq:Jacobi-plus}
 \vartheta_+(x)=\frac12\sqrt{\frac\pi x}-\frac12
 +\sqrt{\frac\pi x}\,\vartheta_+\!\left(\frac{\pi^2}{x}\right).
\end{equation}
Periodicity of $P_K$ and the elementary terms in \cref{eq:MK} give
\begin{equation}\label{eq:MK-difference}
 M_K(x)-M_K(4x)=\frac12\sqrt{\frac\pi x}-\frac12.
\end{equation}
Thus $D(x)=K(x)-M_K(x)$ satisfies
\begin{equation}\label{eq:D-difference}
 D(x)-D(4x)=\sqrt{\frac\pi x}\,
 \vartheta_+\!\left(\frac{\pi^2}{x}\right).
\end{equation}
Replace $x$ by $x/4$, solve for $D(x)$, and iterate toward zero.  The Mellin
flatness $D(t)=\bigO_N(t^N)$ makes the terminal term vanish and yields
\[
 D(x)=-\sqrt{\frac\pi x}
 \sum_{r=1}^{\infty}2^r
 \vartheta_+\!\left(\frac{4^r\pi^2}{x}\right),
\]
which is \cref{eq:theta-completion}.  The first term is $(r,n)=(1,1)$.
The next possible exponent is $16\pi^2/x$, twelve units of $\pi^2/x$ beyond
the leading $4\pi^2/x$, proving \cref{eq:K-flat-leading}.
\end{proof}

\begin{remark}[Size of the periodic ripple]
The first nonzero Fourier coefficient of $P_K$ is
\[
 c_1=\frac{\Gamma(\pi\ii/L_2)\zeta(2\pi\ii/L_2)}{2L_2},
 \qquad |c_1|=0.0011153475185173782067\ldots.
\]
Thus the first-harmonic peak-to-center amplitude is
$2|c_1|=0.0022306950370347564135\ldots$.  The ripple is small but not
formal noise; it is an exact complex-dimension contribution.
\end{remark}

\begin{table}[ht]
\centering
\caption{The exact flat correction $K(x)-M_K(x)$ and the theta completion.}
\label{tab:theta-check}
\begin{tabular}{@{}rcc@{}}
\toprule
$x$ & direct difference & theta-series value\\
\midrule
$1$ & $-2.5371492292\times10^{-17}$ & $-2.5371492292\times10^{-17}$\\
$2$ & $-6.7059525212\times10^{-9}$  & $-6.7059525212\times10^{-9}$\\
$5$ & $-5.9029369697\times10^{-4}$  & $-5.9029369697\times10^{-4}$\\
\bottomrule
\end{tabular}
\end{table}

\section{Endpoint tilting: truncated exponentials, mod-Gamma limits, and Lambert $W$}

\subsection{Exact negative-Laplace decomposition}

Let
\begin{equation}\label{eq:L-g-def}
 L(s)=\E\e^{-sX},\qquad g(s)=\log L(s),\qquad s>0.
\end{equation}
The binary random series gives
\begin{equation}\label{eq:L-product}
 L(s)=\prod_{k\ge1}\frac{1-\e^{-s/2^k}}{s/2^k}.
\end{equation}
The current exact endpoint theory supplies a real-analytic one-periodic
function $R$ and an exponentially small positive error $E$ such that
\begin{equation}\label{eq:g-decomp}
 \boxed{
 g(s)=-\frac{(\log s)^2}{2\log2}+\frac{\log s}{2}
 +R(\log_2s)+E(s),}
\end{equation}
with
\begin{equation}\label{eq:E-bound}
 0\le E(s)\le\frac{\e^{-s}}{(1-\e^{-s})^2}.
\end{equation}
Termwise differentiation of the defining forward-tail series gives the
corresponding exponentially small bounds for each fixed derivative of $E$.

\subsection{The tilted binary coordinates}

Define the exponentially tilted probability measure
\begin{equation}\label{eq:tilt-def}
 \Prob_s(\dd x)=\frac{\e^{-sx}}{L(s)}\Prob(X\in\dd x).
\end{equation}
Let
\begin{equation}\label{eq:lambda-W}
 \lambda_k=s2^{-k},\qquad W_{k,s}=\lambda_kV_k.
\end{equation}

\begin{theorem}[Exact tilted-coordinate law]\label{thm:tilted-coordinates}
Under $\Prob_s$, the coordinates $V_k$ remain independent.  The variables
$W_{k,s}$ have densities
\begin{equation}\label{eq:truncated-exp-density}
 f_{\lambda_k}(w)=
 \frac{\e^{-w}}{1-\e^{-\lambda_k}}
 \bm1_{(0,\lambda_k)}(w),
\end{equation}
and
\begin{equation}\label{eq:sX-sumW}
 sX=\sum_{k\ge1}W_{k,s}.
\end{equation}
\end{theorem}

\begin{proof}
The Radon--Nikodym factor factorizes:
\[
 \e^{-sX}=\prod_{k\ge1}\e^{-\lambda_kV_k}.
\]
Normalizing each uniform coordinate gives density
$\lambda_k\e^{-\lambda_kv}/(1-\e^{-\lambda_k})$ on $0<v<1$.
The change of variable $w=\lambda_kv$ yields
\cref{eq:truncated-exp-density}, and summing gives \cref{eq:sX-sumW}.
\end{proof}

This theorem is the probabilistic mechanism behind the endpoint saddle.  About
$\log_2s$ coordinates have $\lambda_k\gtrsim1$ and behave almost as independent
unit exponentials; the remaining coordinates form a uniformly bounded tail.

\subsection{Tilted cumulants and exact phase dependence}

Let
\begin{equation}\label{eq:K-r-def}
 K_r(s)=\kappa_r^{(s)}(sX)=(-1)^rs^rg^{(r)}(s),
\end{equation}
where the superscript indicates cumulants under $\Prob_s$.  Put
\begin{equation}\label{eq:alpha-beta}
 \alpha(s)=K_1(s)=-sg'(s),\qquad
 \beta(s)=K_2(s)=s^2g''(s).
\end{equation}
Writing $\theta=\log_2s$ and differentiating \cref{eq:g-decomp} gives
\begin{align}
 \alpha(s)&=\log_2s-\frac12-
 \frac{R'(\theta)}{\log2}+\bigO(s\e^{-s}),
 \label{eq:alpha-asymp}\\
 \beta(s)&=\alpha(s)-\frac1{\log2}
 +\frac{R''(\theta)}{(\log2)^2}+\bigO(s^2\e^{-s}).
 \label{eq:beta-asymp}
\end{align}
In particular $\beta(s)=\log_2s+\bigO(1)$.  The cumulants obey the exact
recursion
\begin{equation}\label{eq:K-recur}
 K_{r+1}(s)=rK_r(s)-sK_r'(s).
\end{equation}

\begin{theorem}[Phase-resolved mod-Gamma convergence]\label{thm:mod-gamma}
Let $s_j\to\infty$ with fractional parts
$\{\log_2s_j\}\to\theta\in[0,1]$.  For real $t<1$, locally uniformly in $t$,
\begin{equation}\label{eq:mod-gamma}
 \boxed{
 \E_{s_j}\e^{t s_jX}(1-t)^{\log_2s_j}
 \longrightarrow\Psi_\theta(t),}
\end{equation}
where
\begin{equation}\label{eq:Psi-theta}
 \boxed{
 \Psi_\theta(t)=\exp\!\left[
 -\frac{\log^2(1-t)}{2\log2}+\frac12\log(1-t)
 +R\!\left(\theta+\log_2(1-t)\right)-R(\theta)
 \right].}
\end{equation}
The convergence also holds in a complex neighborhood of the origin.
\end{theorem}

\begin{proof}
The tilted moment generating function is exact:
\begin{equation}\label{eq:tilted-mgf-ratio}
 \E_s\e^{tsX}=\frac{L((1-t)s)}{L(s)}.
\end{equation}
Insert \cref{eq:g-decomp}.  If $v=\log(1-t)$ and $u=\log s$, the quadratic
part contributes
\[
 -\frac{(u+v)^2-u^2}{2\log2}+\frac v2
 =-\frac{uv}{\log2}-\frac{v^2}{2\log2}+\frac v2.
\]
Multiplication by $(1-t)^{\log_2s}$ cancels the term $-uv/\log2$.  Periodicity
and the assumed phase convergence give the $R$-difference in
\cref{eq:Psi-theta}; the two exponentially small errors vanish locally
uniformly.  The same argument works for complex $t$ near zero using the
principal logarithm.
\end{proof}

\begin{corollary}[Phase-resolved cumulants]\label{cor:tilted-cumulants}
For every fixed $r\ge1$,
\begin{equation}\label{eq:Kr-asymp}
 K_r(s)=(r-1)!\log_2s+
 \eta_r(\{\log_2s\})+\smallo(1),
\end{equation}
where $\eta_r$ is an explicit analytic one-periodic function obtained by
$r$ differentiations of $\log\Psi_\theta(t)$ at $t=0$.
\end{corollary}

\subsection{Central and local concentration under the tilt}

\begin{theorem}[Tilted central limit theorem with Berry--Esseen rate]
\label{thm:tilted-CLT}
As $s\to\infty$,
\begin{equation}\label{eq:tilted-CLT}
 \frac{sX-\alpha(s)}{\sqrt{\beta(s)}}
 \Longrightarrow N(0,1).
\end{equation}
Moreover
\begin{equation}\label{eq:Berry-Esseen}
 \sup_{x\in\R}
 \left|\Prob_s\!\left(
 \frac{sX-\alpha(s)}{\sqrt{\beta(s)}}\le x\right)-\Phi_{\mathrm N}(x)\right|
 =\bigO\!\left((\log s)^{-1/2}\right),
\end{equation}
where $\Phi_{\mathrm N}$ is the standard normal distribution function.
\end{theorem}

\begin{proof}
Use the independent triangular array in \cref{thm:tilted-coordinates}.  If
$\lambda\ge1$, a truncated exponential has every fixed centered absolute
moment bounded uniformly in $\lambda$.  There are $\log_2s+\bigO(1)$ such
coordinates.  If $\lambda<1$, the corresponding variable lies in an interval
of length $\lambda$, and the sum of its third centered absolute moments is
bounded by a convergent geometric series.  Therefore
\begin{equation}\label{eq:third-abs-sum}
 \sum_k\E_s|W_{k,s}-\E_sW_{k,s}|^3=\bigO(\log s).
\end{equation}
By \cref{eq:beta-asymp}, the total variance is asymptotic to $\log_2s$.
The Berry--Esseen inequality gives
$\bigO(\log s)/(\log s)^{3/2}$, proving \cref{eq:Berry-Esseen}; the CLT follows.
\end{proof}

\subsection{Lambert-$W_{-1}$ localization}

The typical endpoint location under the tilt is
\begin{equation}\label{eq:xs}
 x_s=\E_sX=\frac{\alpha(s)}s.
\end{equation}
Set
\begin{equation}\label{eq:y-def}
 y=(\log2)sx_s=(\log2)\alpha(s).
\end{equation}

\begin{theorem}[Phase-corrected Lambert localization]\label{thm:Lambert-localization}
As $s\to\infty$,
\begin{equation}\label{eq:y-Lambert-implicit}
 y\e^{-y}=(\log2)x_s
 \exp\!\left(\frac{\log2}{2}+R'(\log_2s)+\smallo(1)\right).
\end{equation}
Equivalently,
\begin{equation}\label{eq:y-W}
 \boxed{
 y=-W_{-1}\!\left(-(
 \log2)x_s\exp\!\left(\frac{\log2}{2}+R'(\log_2s)+\smallo(1)\right)
 \right).}
\end{equation}
After suppressing the bounded phase correction, one obtains the useful
phase-blind equivalent
\begin{equation}\label{eq:s-W-leading}
 s\sim
 -\frac{W_{-1}(-\sqrt2(\log2)x_s)}{(\log2)x_s}.
\end{equation}
\end{theorem}

\begin{proof}
Equation \cref{eq:alpha-asymp} gives
\[
 y=\log s-\frac{\log2}{2}-R'(\log_2s)+\smallo(1).
\]
Since $x_s=(y/\log2)/s$, elimination of $s$ yields
\cref{eq:y-Lambert-implicit}.  The relevant solution has $y\to\infty$, so it
is the $W_{-1}$ branch.  A bounded multiplicative change in the small argument
of $W_{-1}$ changes its value by $\bigO(1)$; division by
$y\asymp\log(1/x_s)$ proves the phase-blind equivalent.
\end{proof}

\begin{remark}
Lambert $W$ is not merely a formal inversion of the leading endpoint
asymptotic.  The tilted law concentrates in a Gaussian window of width
$\sqrt{\beta(s)}/s\asymp\sqrt{\log s}/s$ around $x_s$, and
\cref{thm:Lambert-localization} identifies that actual probabilistic center.
\end{remark}

\section{A universal endpoint-transfer theorem}

\subsection{Universal correction polynomials}

For fixed $c>0$ and an integer $N\ge1$, define the one-sided endpoint moment
\begin{equation}\label{eq:MNc}
 M_N(c)=\E\left[(1-cX)^N\bm1_{\{X\le1/c\}}\right].
\end{equation}
Set $s=cN$ and $Z=sX$ under the tilted law $\Prob_s$.  Then
\begin{equation}\label{eq:MN-tilt-exact}
 \frac{M_N(c)}{L(cN)}
 =\E_{cN}\left[\e^Z\left(1-\frac ZN\right)^N
 \bm1_{\{Z\le N\}}\right].
\end{equation}
Define universal polynomials $P_r$ by the formal identity
\begin{equation}\label{eq:P-polynomials-def}
 \e^z(1-zt)^{1/t}
 =\exp\!\left(-\sum_{j\ge2}\frac{z^jt^{j-1}}j\right)
 =\sum_{r\ge0}P_r(z)t^r.
\end{equation}
The first four are
\begin{align}
 P_0(z)&=1,\nonumber\\
 P_1(z)&=-\frac{z^2}{2},\nonumber\\
 P_2(z)&=\frac{z^4}{8}-\frac{z^3}{3},\label{eq:P-first}\\
 P_3(z)&=-\frac{z^6}{48}+\frac{z^5}{6}-\frac{z^4}{4}.\nonumber
\end{align}

\begin{theorem}[Universal endpoint transfer]\label{thm:endpoint-transfer}
Fix $c>0$ and an integer $R\ge0$.  As $N\to\infty$,
\begin{equation}\label{eq:endpoint-transfer}
 \boxed{
 \frac{M_N(c)}{L(cN)}=
 \sum_{r=0}^{R}\frac{\E_{cN}P_r(Z)}{N^r}
 +\bigO_{R,c}\!\left(
 \frac{(\log N)^{2R+2}}{N^{R+1}}\right),
 \qquad Z=cNX.}
\end{equation}
Every coefficient is an explicit polynomial in the tilted cumulants
$K_1(cN),\dots,K_{2r}(cN)$ and hence in the periodic endpoint function $R$ and
its derivatives.
\end{theorem}

\begin{proof}
Equation \cref{eq:MN-tilt-exact} is obtained by multiplying and dividing the
integrand by $\e^{-cNX}$.  On $0\le Z\le(\log N)^2$, Taylor's theorem applied
to \cref{eq:P-polynomials-def} gives a remainder bounded by
\[
 C_{R,c}N^{-R-1}(1+Z)^{2R+2}.
\]
The mod-Gamma estimate, or directly the independent truncated-exponential
representation, implies
\begin{equation}\label{eq:tilted-moment-bound}
 \E_sZ^k=\bigO_k((\log s)^k)
\end{equation}
for each fixed $k$.  This gives the displayed remainder on the central region.

For the complement, use Markov's inequality with any fixed $0<t<1$ and
\cref{eq:tilted-mgf-ratio}.  The mod-Gamma theorem gives
$\E_s\e^{tZ}=\bigO_t(s^{-\log_2(1-t)})$, only a fixed power of $s$, while
$\exp(-t(\log N)^2)$ beats every power of $N$.  The tail, including the region
$Z>N$, is therefore smaller than the asserted error.  Finally, raw moments are
complete Bell polynomials in cumulants, which proves the last assertion.
\end{proof}

\paragraph{Lean status.}
The finite algebraic last sentence is now formal.  Under an arbitrary
probability measure, with integrability assumed only for the powers occurring
in a chosen polynomial, the theorem
\texttt{integral\_eval}\textsubscript{2}\allowbreak
\texttt{\_eq\_sum\_completeBell}\allowbreak
\texttt{\_momentCumulant} in
\path{PolynomialExpectationCumulant.lean} rewrites its expectation as the
coefficient-weighted complete Bell transform of the formal cumulants of its
raw moments.  The universal polynomials \(P_r\), their formal exponential
identity, recurrence, scalar evaluation, and first values are developed in
\path{EndpointTransferPolynomials.lean}.  These finite identities do
\emph{not} identify the formal cumulants with the analytic tilted cumulants
\(K_r(s)\), nor do they formalize the asymptotic expansion
\eqref{eq:endpoint-transfer}, the mod-Gamma estimate, or its remainder and
tail bounds; those analytic steps remain frontier work.

\subsection{The first two corrections}

Write, at $s=cN$,
\begin{equation}\label{eq:abcd}
 a=K_1(s)=\alpha(s),\quad b=K_2(s)=\beta(s),\quad
 c_3=K_3(s),\quad c_4=K_4(s).
\end{equation}
Then
\begin{align}
 \E_sZ^2&=a^2+b,\label{eq:raw2}\\
 \E_sZ^3&=a^3+3ab+c_3,\label{eq:raw3}\\
 \E_sZ^4&=a^4+6a^2b+3b^2+4ac_3+c_4.\label{eq:raw4}
\end{align}
Thus
\begin{equation}\label{eq:transfer-second}
 \frac{M_N(c)}{L(cN)}=
 1-\frac{\E_sZ^2}{2N}
 +\frac{\E_sZ^4/8-\E_sZ^3/3}{N^2}
 +\bigO\!\left(\frac{(\log N)^6}{N^3}\right).
\end{equation}

\subsection{Ordinary Fabius moments}

By symmetry $X\dist1-X$,
\begin{equation}\label{eq:mn-M}
 m_n=\E X^n=\E(1-X)^n=M_n(1).
\end{equation}

\begin{corollary}[Large ordinary moments]\label{cor:ordinary-moments}
With all tilted quantities evaluated at $s=n$,
\begin{equation}\label{eq:mn-second}
 \boxed{
 \frac{m_n}{L(n)}=
 1-\frac{\alpha(n)^2+\beta(n)}{2n}
 +\frac{\E_nZ^4/8-\E_nZ^3/3}{n^2}
 +\bigO\!\left(\frac{(\log n)^6}{n^3}\right).}
\end{equation}
In logarithmic form,
\begin{equation}\label{eq:logmn-first}
 \log m_n=g(n)-\frac{\alpha(n)^2+\beta(n)}{2n}
 +\bigO\!\left(\frac{(\log n)^4}{n^2}\right).
\end{equation}
If $u=\log n$, $\theta=u/\log2$, and
\begin{equation}\label{eq:An}
 A_n=\frac{u}{\log2}-\frac12-\frac{R'(\theta)}{\log2},
\end{equation}
then
\begin{align}
 \log m_n={}&-\frac{u^2}{2\log2}+\frac u2+R(\theta)\nonumber\\
 &-\frac1{2n}\left[
 A_n^2+A_n-\frac1{\log2}+\frac{R''(\theta)}{(\log2)^2}
 \right]
 +\bigO\!\left(\frac{u^4}{n^2}\right).
 \label{eq:logmn-R}
\end{align}
\end{corollary}

\subsection{Centered moments and the $q$/MZV coefficient bridge}

Symmetry about $1/2$ gives
\begin{equation}\label{eq:mu-M}
 \mu_{2r}=\E(1-2X)^{2r}=2M_{2r}(2).
\end{equation}
Here $N=2r$ and the Laplace scale is $s=cN=4r$.

\begin{corollary}[Large centered moments]\label{cor:centered-moments}
With the tilted quantities evaluated at $s=4r$,
\begin{align}
 \frac{\mu_{2r}}{2L(4r)}={}&
 1-\frac{\alpha(4r)^2+\beta(4r)}{4r}\nonumber\\
 &+\frac{\E_{4r}Z^4/8-\E_{4r}Z^3/3}{4r^2}
 +\bigO\!\left(\frac{(\log r)^6}{r^3}\right).
 \label{eq:mu-transfer}
\end{align}
Consequently the exact coefficient $C_r$ in the $q$-binomial/MZV expansion
\cref{eq:C-partition} satisfies
\begin{equation}\label{eq:Cr-asymp}
 \boxed{\begin{aligned}
 \log C_r={}&2r\log(2\pi)-\log(2r)!+\log2+g(4r)\\
 &-\frac{\alpha(4r)^2+\beta(4r)}{4r}
 +\bigO\!\left(\frac{(\log r)^4}{r^2}\right).
 \end{aligned}}
\end{equation}
Since $R$ has period one,
$R(\log_2(4r))=R(\log_2r)$: the same binary phase controls both the endpoint
Laplace law and the large-order $q$/MZV coefficients.
\end{corollary}

\begin{proof}
Apply \cref{thm:endpoint-transfer} to \cref{eq:mu-M} and then use
\cref{eq:C-moment}.  Taking the logarithm changes the second-order remainder
only by $\bigO((\log r)^4/r^2)$.
\end{proof}

\begin{boxedremark}
Equation \cref{eq:Cr-asymp} is the main bridge result.  The left side is an
exact coefficient assembled from Gaussian $q$-binomials and MZVs.  The right
side is governed by the endpoint Laplace product, its log-periodic correction,
and the Lambert-$W_{-1}$ localization of the tilted mass.  Neither side is an
approximate restatement of the other; the transfer theorem proves why they
must agree at large order and supplies computable corrections.
\end{boxedremark}

\subsection{Numerical validation}

The ordinary moments can be generated exactly by
\begin{equation}\label{eq:mn-recurrence}
 (2^n-1)m_n=\sum_{k=1}^{n}\binom nk\frac{m_{n-k}}{k+1},
 \qquad m_0=1.
\end{equation}
Table \ref{tab:ordinary-transfer} compares the exact ratio $m_n/L(n)$ with the
one- and two-correction truncations of \cref{eq:mn-second}.

\begin{table}[ht]
\centering
\caption{Endpoint transfer for ordinary moments.}
\label{tab:ordinary-transfer}
\begin{tabular}{@{}rccc@{}}
\toprule
$n$ & exact ratio & one correction & two corrections\\
\midrule
20   & 0.650564349158 & 0.575338248846 & 0.658507245134\\
50   & 0.740712148665 & 0.698394360314 & 0.745157800500\\
100  & 0.809797566277 & 0.787758642159 & 0.811530920395\\
200  & 0.868346118425 & 0.858160052080 & 0.868893882715\\
500  & 0.924530320090 & 0.921307792046 & 0.924627286470\\
1000 & 0.952442261825 & 0.951188127893 & 0.952465665513\\
\bottomrule
\end{tabular}
\end{table}

For centered moments the natural comparison is
$\mu_{2r}/(2L(4r))$.

\begin{table}[ht]
\centering
\caption{Endpoint transfer for centered even moments.}
\label{tab:centered-transfer}
\begin{tabular}{@{}rccc@{}}
\toprule
$r$ & exact ratio & one correction & two corrections\\
\midrule
10  & 0.516638138954 & 0.334240960496 & 0.558725000470\\
20  & 0.623451024335 & 0.521571821828 & 0.641626320479\\
50  & 0.755044449961 & 0.716320104160 & 0.759255426702\\
100 & 0.834124742928 & 0.817440783081 & 0.835296570356\\
\bottomrule
\end{tabular}
\end{table}

The second correction is already visibly asymptotic rather than merely
qualitative.  Its slower start for centered moments is expected: their
relevant Laplace scale is $4r$, but the polynomial exponent is only $2r$, so
the first correction carries the larger factor
$(\alpha(4r)^2+\beta(4r))/(4r)$.

\section{Synthesis: one determinant, two probability laws, three asymptotic regimes}

The results can be summarized by the single determinant
\begin{equation}\label{eq:synthesis-A}
 \mathcal A(w)=
 \underbrace{\Phi(\ii\sqrt w)}_{\text{sinc/zero geometry}}
 =\underbrace{\E\exp(2\pi\sqrt w\,Y)}_{\text{compact law}}
 =\underbrace{\bigl(\E\e^{-wV_1}\bigr)^{-1}}_{\text{gamma-convolution dual}}.
\end{equation}
Each equality exposes a different rigidity.

\begin{enumerate}
\item The sinc product gives integer zeros, dyadic multiplicities, and the
      Thue--Morse sign of real Fourier lobes.
\item The compact-law MGF gives exact centered moments, optimal sub-Gaussian
      concentration, and endpoint tilting.
\item The reciprocal Laplace transform gives a GGC, Thorin measure, complete
      Bernstein structure, and the L\'evy heat kernel.
\item Expanding by dyadic layers produces Gaussian $q$-binomial occupancies;
      expanding by distinct integer sites produces symmetrized MZVs.
\item Mellin analysis of the spectral heat kernel and of the negative-Laplace
      product produces log-periodic phases.  Endpoint transfer proves that
      those phases survive into large moment order.
\end{enumerate}

This synthesis also clarifies a striking duality.  The compact Fabius--Rvachev
law is strictly variance-optimal sub-Gaussian but not infinitely divisible.
The reciprocal Wick rotation is the Laplace transform of an explicitly
infinitely divisible GGC.  Their cumulants are related by the alternating
formula \cref{eq:cumulant-duality}; their sign structure is related by Thorin
mass parity; and their asymptotics are tied by the same dyadic complex
dimensions.

\section{Next frontier targets}

\subsection{Phase-resolved Edgeworth and local limits}

The proof of \cref{thm:tilted-CLT} contains more information than the stated
Berry--Esseen bound.  The independent summands have uniformly controlled
moments, the variance diverges logarithmically, and
\cref{cor:tilted-cumulants} gives every cumulant with an explicit periodic
constant term.  This strongly suggests the following sharpened statement.

\begin{conjecture}[Phase-resolved local Edgeworth expansion]
\label{conj:Edgeworth}
Along every sequence $\{\log_2s_j\}\to\theta$, the tilted density of
$(s_jX-\alpha(s_j))/\sqrt{\beta(s_j)}$ admits a full local Edgeworth expansion
in powers of $(\log s_j)^{-1/2}$, uniformly on windows
$|x|\le(\log s_j)^{1/6-\varepsilon}$.  Every coefficient is a polynomial in
$x$ whose constants are analytic one-periodic functions of $\theta$.
\end{conjecture}

A proof should combine the explicit characteristic function
$L((1-\ii t/s)s)/L(s)$ with a triangular-array Cram\'er bound.  The compact
support of each digit and the presence of many near-exponential coordinates
make the nonlattice estimate plausible, but it should be checked uniformly in
the binary phase.

\subsection{All-orders moment transfer}

Theorem \ref{thm:endpoint-transfer} is already all-orders in principle.  Two
concrete tasks remain:
\begin{enumerate}
\item generate $P_r$ by Bell polynomials and express
      $\E_sP_r(Z)$ directly in $R,R',\dots,R^{(2r)}$;
\item combine the result with Stirling's expansion in \cref{eq:Cr-asymp} to
      produce a closed transseries for $C_r$, including the periodic phase and
      exponentially small endpoint terms.
\end{enumerate}
This would turn the q/MZV coefficient sequence into a precise asymptotic probe
of the negative-Laplace periodic function.

\subsection{Inverse recovery of the periodic endpoint function}

Equation \cref{eq:Cr-asymp} can be read in reverse.  Exact rational moments
produce exact $C_r$; after subtracting the explicit Stirling and quadratic-log
terms, the residual samples $R(\log_2r)$ with a controlled error.  Since
$\{\log_2r\}$ is dense modulo one, this suggests a certified reconstruction
scheme for $R$ from exact arithmetic alone.  The correction terms involving
$R'$ and $R''$ make the inverse problem nonlinear but also provide derivative
information.  A practical scheme would alternate Fourier reconstruction of
$R$ with the transfer correction until the residual falls below the exact
moment error bound.

\subsection{A $q$-deformed gamma dual}

The corpus defines symmetric $q$-integer deformations of the sinc determinant.
Whenever the deformed zeros remain positive after Wick rotation, the same
argument as \cref{thm:gamma-dual} constructs a discrete Thorin measure at the
squared $q$-integers.  The main questions are:
\begin{itemize}
\item whether the deformed determinant remains type-I Laguerre--P\'olya for the
      full parameter range;
\item whether its Thorin mass admits a deformed digit-parity interpretation;
\item how the $q\uparrow1$ and $q\downarrow0$ limits interact with the
      perpetuity and heat-kernel complex dimensions.
\end{itemize}
These questions are exact and experimentally accessible; no conjectural
analytic continuation is needed to begin.

\section{Lean formalization roadmap}

The new results split naturally into modules with very different library
requirements.

\begin{boxedremark}
\textbf{Post-snapshot algebraic progress.}
The later modules \texttt{ExponentialPartition.lean},
\texttt{ExponentialBell.lean}, and \texttt{MomentCumulantAlgebra.lean}
now formalize, over every commutative
\(\mathbb Q\)-algebra, the universal weighted-partition formula for an
exponential coefficient and identify it with the recursive
\texttt{SaddleExpansion.expCoeff} engine, then factorially normalize it
into named complete-Bell and inverse moment--cumulant transforms with the
classical Bell recurrence and normalized inverse laws.  This supplies reusable
algebra for the Bell-polynomial sentence in
\cref{thm:endpoint-transfer} and for the coefficient stage of the
proposed \texttt{MomentTransfer.lean}.  It does not establish the
tilted-cumulant formulas, their probabilistic interpretation, the
Taylor and tail estimates, or the hierarchy-specific endpoint identities.
Consequently none of the status labels in this report
changes.
\end{boxedremark}

\subsection{Low-risk algebraic modules}

\begin{enumerate}
\item \texttt{QBinomialMZVMoments.lean}: prove the finite-depth occupancy
      expansion over finite spectral cutoffs, then pass first in the spectral
      cutoff and then in depth.  The partition-lattice identity can be stated
      initially as a finite inclusion--exclusion theorem, avoiding a general
      MZV library.
\item \texttt{SpectralGammaCumulants.lean}: formalize the algebraic identities
      for Laplace exponents and cumulants before constructing the countable
      probability space.  The equality
      $\kappa_r=\tau(r-1)!\zeta(2r)/(1-4^{-r})$ is mostly existing Dirichlet
      series infrastructure.
\item \texttt{ThorinParity.lean}: combine the already formalized lobe-sign
      theorem with Legendre's digit-sum formula and a finite cumulative mass
      definition.  This should be short.
\end{enumerate}

\subsection{Analytic modules}

\begin{enumerate}
\item \texttt{GammaConvolutionDual.lean}: construct partial sums of independent
      gamma variables, prove $L^1$ convergence from the expectation sum, and
      identify the Laplace transform by dominated convergence.
\item \texttt{LevyHeatThetaCompletion.lean}: establish
      $K(x)-K(4x)=\vartheta_+(x)$, import or prove the Jacobi theta transform,
      and telescope the flat correction.  The exact identity
      \cref{eq:theta-completion} is more formalization-friendly than a contour
      shift because all sums are positive or exponentially dominated.
\item \texttt{EndpointTilt.lean}: prove coordinate factorization under the
      tilted finite product, pass to the infinite law, and derive the cumulant
      recursion.
\item \texttt{MomentTransfer.lean}: define $P_r$ through formal power series,
      prove a finite-order Taylor bound on a logarithmic window, and control the
      complement with the exact Laplace ratio.
\end{enumerate}

\subsection{Library-heavy modules}

The Laguerre--P\'olya and PF$_\infty$ statements may require new general
infrastructure.  A pragmatic first formal target is finite truncation:
\begin{itemize}
\item prove real nonpositive roots of each finite polynomial;
\item prove total positivity of its coefficient Toeplitz matrix by repeated
      convolution with $(1,\alpha)$;
\item derive every fixed finite minor inequality for the limit by coefficient
      convergence.
\end{itemize}
This avoids formalizing the full Aissen--Schoenberg--Whitney theorem while
still obtaining any requested moment inequality.

\appendix

\section{Exact coefficients and normalized Tur\'an ratios}

For reference, the first five nontrivial centered moments are
\begin{equation}\label{eq:moment-table-display}
\begin{array}{c|c}
 r&\mu_{2r}\\
\hline
1&1/9\\
2&19/675\\
3&583/59535\\
4&132809/32531625\\
5&46840699/24405225075
\end{array}
\end{equation}
The normalized Tur\'an ratio
\begin{equation}\label{eq:Turan-ratio}
 T_r=\frac{2r+1}{2r-1}
 \frac{\mu_{2r}^2}{\mu_{2r-2}\mu_{2r+2}}
 =\frac{\gamma_r^2}{\gamma_{r-1}\gamma_{r+1}}
\end{equation}
is at least one.  Its first values are
\begin{equation}\label{eq:Turan-values}
 T_1=\frac{25}{19}=1.315789\ldots,
 \qquad
 T_2=\frac{17689}{14575}=1.21365\ldots,
 \qquad
 T_3=1.16829\ldots.
\end{equation}
The approach toward one is consistent with the large-order endpoint
asymptotic: the Jensen coefficients become asymptotically close to a geometric
sequence on the scale of one index, while retaining strict log-concavity.

\section{Proof-dependency map}

\begin{longtable}{@{}p{0.22\textwidth}p{0.70\textwidth}@{}}
\toprule
Result & Minimal dependencies\\
\midrule
\endhead
$q$/MZV formula & finite $q$-binomial theorem; normal convergence; set-partition
M\"obius inversion\\
Laguerre--P\'olya & canonical product convergence; finite real-rooted
truncations\\
PF$_\infty$ & Aissen--Schoenberg--Whitney characterization\\
Jensen hierarchy & Jensen characterization of the Laguerre--P\'olya class\\
Sub-Gaussianity & $\log(1+x)<x$; exact spectral sum $p_1=2\pi^2/9$\\
Gamma dual & gamma Laplace transform; Tonelli; regrouping equal rates\\
Thorin parity & $\sum_{m\le N}(1+v_2(m))=2N-s_2(N)$; existing lobe sign\\
Theta completion & dyadic heat equation; Jacobi theta transform; telescoping\\
Mod-Gamma limit & exact negative-Laplace decomposition; periodicity of $R$\\
Berry--Esseen & tilted independent digits; third-moment bound; variance asymptotic\\
Moment transfer & exponential tilt; universal Taylor polynomials; tilted moment
and tail bounds\\
\bottomrule
\end{longtable}

\section{Corpus map used in the audit}

The audit followed the repository's revision structure rather than reading the
directory as a flat bag of papers.  The decisive source groups were:
\begin{description}[style=nextline,leftmargin=2.2em,labelindent=0pt]
\item[Current expository core.]
\texttt{Fabius\_Function\_and\_Rvachev\_Up}, the current Lean walkthrough,
the mathematical glossary, and the non-elementarity article.

\item[Archived foundational studies.]
The earlier primary expositions, both missing-parts supplements, the two
$q$-special-function expositions, the small-argument and inverse-function
studies, and the early notes on Fabius asymptotics and iterated Thue--Morse
summation.

\item[Exact arithmetic and convergence studies.]
The dyadic formula/alternative-representation papers, the two
dyadic-to-asymptotic bridges, the Thue--Morse formula atlases, both
Thue--Morse convergence-rate reports, and both repeated-integration reports.

\item[Live frontier synthesis.]
The consolidated \texttt{non-formalized-research-frontiers} manuscript, the two
Fabius--Rvachev frontier reports, and the spectral-determinant/$q$-deformation
results bundle.

\item[Fourier-decay campaign.]
The original sinc-decay question, all numbered response waves, the final
syntheses, gentle guide, comparative audit, and second-wave audit.  The latest
second-wave audit was treated as authoritative where older agents disagreed.

\item[Reference transcriptions.]
The corrected and source \LaTeX{} transcriptions under \texttt{docs/papers},
including the work on the Fabius function, arithmetic properties, and lacunary
products.  They were used for historical and literature context, not as a
substitute for checking the repository's current normalization.
\end{description}

\section{Bibliographic notes}

\begin{thebibliography}{99}

\bibitem{ASW1952}
M.~Aissen, I.~J.~Schoenberg, and A.~Whitney,
\emph{On the generating functions of totally positive sequences. I},
Journal d'Analyse Math\'ematique \textbf{2} (1952), 93--103.

\bibitem{AndrewsAskeyRoy}
G.~E.~Andrews, R.~Askey, and R.~Roy,
\emph{Special Functions}, Cambridge University Press, 1999.

\bibitem{Bondesson}
L.~Bondesson,
\emph{Generalized Gamma Convolutions and Related Classes of Distributions and
Densities}, Lecture Notes in Statistics 76, Springer, 1992.

\bibitem{Corless}
R.~M.~Corless, G.~H.~Gonnet, D.~E.~G.~Hare, D.~J.~Jeffrey, and D.~E.~Knuth,
\emph{On the Lambert W function}, Advances in Computational Mathematics
\textbf{5} (1996), 329--359.

\bibitem{FlajoletMellin}
P.~Flajolet, X.~Gourdon, and P.~Dumas,
\emph{Mellin transforms and asymptotics: harmonic sums},
Theoretical Computer Science \textbf{144} (1995), 3--58.

\bibitem{GasperRahman}
G.~Gasper and M.~Rahman,
\emph{Basic Hypergeometric Series}, second edition, Cambridge University Press,
2004.

\bibitem{Hoffman}
M.~E.~Hoffman,
\emph{Multiple harmonic series}, Pacific Journal of Mathematics
\textbf{152} (1992), 275--290.

\bibitem{Petrov}
V.~V.~Petrov,
\emph{Limit Theorems of Probability Theory}, Oxford University Press, 1995.

\bibitem{Sato}
K.-I.~Sato,
\emph{L\'evy Processes and Infinitely Divisible Distributions},
Cambridge University Press, 1999.

\bibitem{Schilling}
R.~L.~Schilling, R.~Song, and Z.~Vondra\v{c}ek,
\emph{Bernstein Functions: Theory and Applications}, second edition,
de Gruyter, 2012.

\bibitem{ProveItPrimary}
V.~Reshetnikov,
\emph{The Fabius Function and Rvachev's Up-Function: Exact Arithmetic,
Probability, Fourier Analysis, and Corrected Sharp Asymptotics},
\texttt{ProveIt}, \nolinkurl{Analysis/FabiusFunction/docs/}, current snapshot.

\bibitem{ProveItFrontier}
V.~Reshetnikov,
\emph{Non-Formalized Research Frontiers for the Fabius Function},
\texttt{ProveIt}, \nolinkurl{Analysis/FabiusFunction/docs/non-formalized-research-frontiers/},
current snapshot.

\end{thebibliography}

\end{document}
